
Reinforcement learning from execution feedback (RLEF) has emerged as a paradigm for improving code generation models by providing reward signals derived from program execution outcomes. Prior work has explored various feedback granularities: binary pass/fail rewards \citep{le2022coderl}, continuous test pass rates \citep{liu2023rltf}, and multi-signal combinations incorporating error type information \citep{shojaee2023ppocoder}. Information theory suggests that higher bandwidth signals---those providing more information per gradient update---should enable more precise credit assignment and faster convergence.

This work aimed to provide a controlled comparison of feedback granularities using an ''information bandwidth'' framework. The hypothesis was that under fixed compute budget, higher bandwidth rewards would accelerate convergence to a pass@1 threshold. However, the proof-of-concept scale experiment was insufficient to produce any learning signal across all conditions.

The result was that at proof-of-concept scale (1 epoch, 3 seeds per condition), all nine training runs produced 0.0\% pass@1 on the evaluation set. No condition exhibited learning. Statistical comparison was not possible because all values were at floor with zero variance.

This null result does not falsify the information bandwidth hypothesis. It is a finding about the specific experimental setup: the configuration (REINFORCE algorithm, 1 epoch, CodeLlama-7B-Instruct, MBPP) does not produce measurable learning within the allocated compute budget.

This report documents: (1) the experimental configuration that failed to produce learning, (2) the methodology for future adequate-scale replication, and (3) a diagnostic gap---the absence of training curves prevents distinguishing between insufficient scale, catastrophic forgetting, and implementation error.
